\documentclass[10pt,a4paper]{amsart}
\usepackage[margin=19mm]{geometry}
\usepackage{amsmath,amssymb,mathtools}
\usepackage[scr=rsfs,cal=euler]{mathalfa}
\usepackage{xfrac}
\usepackage[unicode,hidelinks,bookmarks=false]{hyperref}
\newcommand{\ie}{i.e.\ }
\newcommand{\cf}{cf.\ }
\newcommand{\resp}{respectively\ }
\newcommand{\Subsection}{\subsection}
\providecommand{\nobreakdash}{\nobreak\hbox{-}}
\setlength{\emergencystretch}{2em}
\multlinegap0pt
\newtheorem{lem}{Lemma}
\title{Chromatic numbers from edge ideals:\\ Graph classes with vanishing syzygies are polynomially $\chi$-bounded}
\author{Alexander Engstr\"om}
\date{Source: arXiv:2512.21800v2 (CC BY 4.0)}
\begin{document}
\maketitle

\noindent\emph{Source and licence.} Chromatic numbers from edge ideals:\\ Graph classes with vanishing syzygies are polynomially $\chi$-bounded. Version arXiv:2512.21800v2 (CC BY 4.0), distributed by arXiv under the Creative Commons Attribution 4.0 International licence, http://creativecommons.org/licenses/by/4.0/, as recorded in the arXiv licence metadata for this exact version and shown on the abstract page https://arxiv.org/abs/2512.21800v2. Full licence notice: \texttt{licenses/arxiv-2512.21800-CC-BY-4.0.txt}.\par\smallskip

\noindent\textit{Editorial scope.} Counted unit: the article's lem K2extend (source0.tex lines 74--79). The article's complete original proof of it is delivered with it (source0.tex lines 81--132, one \texttt{proof} environment of 52 lines). No mathematics is written, repaired or re-derived: the statement and the proof are the article's own lines, in the article's own notation, and no retained line is substituted or re-flowed.  No further source-local context is needed: every cross-reference of every retained line resolves inside the two delivered spans.  Nothing was omitted from any retained span, so the package carries no omission record.  No retained line contains a citation, so the wrapper carries no bibliography and no bibliography entry.  No retained line contains a figure environment, an image-inclusion command or drawing code, so no image is delivered and the package declares no asset dependency.  Editorial formatting only, in addition: an AMS \texttt{amsart} preamble that restates the article's own declarations and macros used by the retained lines, namely the environments \texttt{lem} on separate counters, together with the standard packages those lines need.  The retained environments print in this document as Lemma 1 in order of appearance, whereas the article prints the same items with the numbering of its own full document.  The complete native source of the article as distributed in the arXiv e-print for this exact version is bundled unchanged as \texttt{sources/191687/source0.tex}.
\begin{lem}\label{K2extend}
If ${\bf H}$ is a set of graphs and the class of ${\bf H}$--free graphs is  $\chi$-bounded by $f_{\bf H},$ then the class of  $K_2 \cup {\bf H  } $--free graphs is $\chi$-bounded by
\[
f_{K_2 \cup {\bf H  } }(\omega) = \sum_{k=1}^{\omega}  (\omega-k+1)  f_{{\bf H  } }(k).
\]
\end{lem}

\begin{proof}
Let $G$ be a $K_2 \cup {\bf H  } $--free graph with a clique on the vertices $1,2,\ldots, \omega=\omega(G).$
For $1\leq i < j \leq \omega$ define
\[
C_{i,j}' = \{ v\in V(G) \setminus \{1,2,\ldots, \omega\} \mid vi,vj \not\in E(G).  \}
\]
For any induced subgraph $H$ of $C_{i,j}'$, the graph $K_2 \cup H$ on $\{i,j\} \cup V(H)$ is an induced subgraph of $G.$ By assumption, $G$ is a $K_2 \cup {\bf H  } $--free graph, and, thus, the induced subgraph of $G$ on $C_{i,j}' $ is an  ${\bf H  } $--free graph.

Set
\[
C_{i,j} = C_{i,j}' \setminus  \left(\,\, \bigcup_{k=1}^{i-1} \bigcup_{l=k+1}^\omega C_{k,l}'   \quad \cup \quad  \bigcup_{l=i+1}^{j-1} C_{i,l}'  \,\,  \right). 
\]
If $v\in C_{i,j}$ and $1\leq k < i$ then $kv \in E(G)$ as 
\[ v\in C_{i,j}' \Rightarrow iv \not\in E(G) \textrm{ and }  jv \not\in E(G) \]
and
\[ v\not\in C_{k,j}' \Rightarrow kv \in E(G) \textrm{ or }  jv \in E(G). \]
If $v\in C_{i,j}$ and $i< l < j$ then $lv \in E(G)$ as 
\[ v\in C_{i,j}' \Rightarrow iv \not\in E(G) \textrm{ and }  jv \not\in E(G) \]
and
\[ v\not\in C_{i,l}' \Rightarrow iv \in E(G) \textrm{ or }  lv \in E(G). \]
It follows that any clique in $G[C_{i,j}]$ extends to a larger clique in $G$ by the $j-2$ vertices
$\{1,2,\ldots,i-1\} \cup \{ i+1,i+2,\ldots, j-1 \}.$ 
Thus
\[ \omega(G[C_{i,j}]) \leq \omega(G)-j+2\]
and
\[ \chi(G[C_{i,j}]) \leq f_{{\bf H  } }( \omega(G)-j+2 ). \]

So far we have partitioned the vertices in $G$ outside the clique on $1,2,\ldots, \omega=\omega(G)$  that has at least two non-edges to the clique. There cannot be a vertex with no non-edges to the clique, as the clique is maximal in $G.$ 
For $1\leq i \leq \omega(G)$ set
\[ D_i'
= \left\{ v\in V(G) \setminus \{1,2,\ldots, \omega\} \mid vi \not\in E(G); \{v1,v2, \ldots, v\omega \} \setminus vi \subseteq E(G) 
%v1,v2,\ldots v(i-1),v(i+1),\ldots v\omega 
 \right \}
\]
If there would be an edge between two vertices $u$ and $v$ in $D_i',$ then the vertices
\[ 1,2, \ldots, i-1, i+1, \ldots, \omega, u, v \]
would constitute a clique on $\omega+1$ vertices, which is a contradiction. Thus, $D_i'$ is independent. As there are no edges from $i$ to vertices of $D_i',$ the set
\[ D_i = \{i\} \cup D_i' \]
is independent for all $i.$ We only need one color for an independent set, and
\[ \chi(G[D_i]) = 1 \leq  f_{{\bf H  } }( 1). \]
Coloring each part of the partition
\[ V(G) \quad = \quad   \bigcup_{i=1}^\omega D_{i}      \quad \cup \quad    \bigcup_{1\leq i < j \leq \omega}  C_{i,j}  \]
gives
\[\begin{array}{rcl} 
\chi(G) & \leq & \displaystyle  \sum_{i=1} ^\omega \chi(G[D_i]) +  \sum_{1\leq i < j \leq \omega}  \chi(G[C_{i,j}]) \\
& \leq & \displaystyle  \sum_{i=1} ^\omega  f_{{\bf H  } }( 1) +  \sum_{1\leq i < j \leq \omega}   f_{{\bf H  } }( \omega-j+2 ) \\
& = & \displaystyle  \omega  f_{{\bf H  } }( 1) +  \sum_{j=2}^\omega   (j-1)f_{{\bf H  } }( \omega-j+2 ) \\
& = & \displaystyle  \sum_{j=2}^{\omega+1}   (j-1)f_{{\bf H  } }( \omega-j+2 ) \\
& = & \displaystyle \sum_{k=1}^{\omega}  (\omega-k+1)  f_{{\bf H  } }(k)
\end{array}
\]
\end{proof}

\end{document}
